\documentclass[11pt,a4paper]{article}
\usepackage[a4paper,margin=2.35cm]{geometry}
% Overleaf-compatible font setup. The document can be compiled with
% pdfLaTeX, XeLaTeX, or LuaLaTeX; no proprietary system fonts are required.
\usepackage{iftex}
\ifPDFTeX
  \usepackage[T1]{fontenc}
  \usepackage[utf8]{inputenc}
  \usepackage{newtxtext}
  \usepackage[scaled=0.92]{helvet}
\else
  \usepackage{fontspec}
  \setmainfont{TeX Gyre Termes}
  \setsansfont{TeX Gyre Heros}[Scale=0.92]
\fi
\usepackage[nopatch=footnote]{microtype}
\usepackage{amsmath,amssymb,mathtools}
\ifPDFTeX
  \usepackage{newtxmath}
\fi
\usepackage{booktabs,longtable,array,tabularx}
\newcolumntype{Y}{>{\RaggedRight\arraybackslash}X}
\usepackage{enumitem}
\usepackage{xcolor}
\usepackage[round,authoryear]{natbib}
\usepackage{hyperref}
\usepackage{fancyhdr}
\usepackage{titlesec}
\usepackage{caption}
\usepackage{float}
\usepackage{hanging}
\usepackage{ragged2e}
\usepackage{setspace}

\definecolor{navy}{HTML}{17324D}
\definecolor{midgray}{HTML}{5F6870}
\titleformat{\section}{\Large\bfseries\color{navy}}{\thesection.}{0.6em}{}
\titleformat{\subsection}{\large\bfseries\color{navy}}{\thesubsection}{0.6em}{}
\titleformat{\subsubsection}{\normalsize\bfseries\color{navy}}{\thesubsubsection}{0.6em}{}
\titlespacing*{\section}{0pt}{1.3em}{0.55em}
\titlespacing*{\subsection}{0pt}{1.0em}{0.4em}
\titlespacing*{\subsubsection}{0pt}{0.8em}{0.3em}
\setlist{leftmargin=1.6em,itemsep=0.3em,topsep=0.4em}
\renewcommand{\arraystretch}{1.2}
\setlength{\parindent}{0pt}
\setlength{\parskip}{0.55em}
\onehalfspacing

\pagestyle{fancy}
\fancyhf{}
\fancyhead[R]{\small\color{midgray}Research Proposal}
\fancyfoot[C]{\thepage}
\setlength{\headheight}{14pt}

\setlength{\emergencystretch}{2em}
\setlength{\bibhang}{1.5em}
\setlength{\bibsep}{0.6em}
\renewcommand{\bibfont}{\small\RaggedRight}
\urlstyle{same}
\Urlmuskip=0mu plus 1mu
\captionsetup[table]{font=small,labelfont=bf,skip=6pt}
\setlength{\LTpre}{0.6em}
\setlength{\LTpost}{0.6em}
\newcommand{\term}[1]{\textit{#1}}
\hypersetup{colorlinks=true,linkcolor=navy,urlcolor=navy,citecolor=navy,pdfauthor={Xiangyu Xu},pdftitle={Delegating Conflict to AI: Escalation and Responsibility Attribution in Intergroup Conflict}}
\fancyhead[L]{\small\color{midgray}AI Decision Agents and Intergroup Conflict}
\newcommand{\AI}{\textit{AI}}
\newcommand{\Baseline}{\textbf{C0 Baseline}}
\newcommand{\AID}{\textbf{C1 AI Delegation}}

\begin{document}

\begin{titlepage}
\centering
\vspace*{2.5cm}
{\sffamily\bfseries\color{navy}\LARGE Research Proposal\par}
\vspace{1.2cm}
{\sffamily\bfseries\Huge Delegating Conflict to AI: Escalation and Responsibility Attribution in Intergroup Conflict\par}
\vspace{1.5cm}
{\sffamily\Large Xiangyu Xu\par}
\vspace{0.8cm}
{\sffamily\small Revised 30 August 2026\par}
\vspace{0.2cm}
{\sffamily\small Research proposal; design and measures remain to be validated\par}
\vfill
\end{titlepage}

\section*{Abstract}

AI decision agents may alter both conflict behavior and subsequent judgments of responsibility. Research on delegation examines whether people transfer decision authority and whether delegation changes behavior. Research on AI responsibility examines how people evaluate humans and AI after an outcome occurs. Research on AI in conflict mainly studies whether AI analysis, advice, or model decisions affect the use of force and conflict escalation. The representative studies reviewed in this proposal cover different stages of the decision process, but they do not jointly explain what happens when a human can voluntarily transfer final authority to AI and real group members subsequently judge responsibility from positions of benefit or harm.

This proposal contains two connected experiments. Study 1 uses a group-based Hawk--Dove game and randomly compares two decision environments. In \Baseline{}, a human must make the final conflict choice. In \AID{}, the human may retain final authority or transfer it to an AI decision agent. Once delegation is confirmed, the AI choice is binding: the human cannot preview, modify, veto, or reconfirm it. Study 1 identifies the overall effect of making AI delegation available on final conflict escalation and separately records actual delegation. Study 2 uses matched decision records. It holds the rules, actions, and outcome constant while randomly varying whether the final choice is recorded as having been made by the human or by AI. This design identifies the effect of the recorded delegation and final-decision arrangement on responsibility ratings.

Where AI is available, responsibility is measured with two independent ratings from 0 to 100: one for the human decision-maker and one for the AI decision agent. The human-only Study 1 baseline collects only the human rating; an absent AI rating is not coded as zero. The ratings are not constrained to sum to 100, allowing both actors to be held responsible at the same time. Study 1 members evaluate their own group's decision-makers after real payoff consequences. Study 2 compares benefiting and harmed records; their connection to evaluators' actual bonuses remains to be specified. The two studies answer separate behavioral and attribution questions. The proposal does not treat post-outcome responsibility ratings as a causal mediator of earlier delegation or escalation.

\textbf{Keywords:} AI delegation; decision agents; conflict escalation; responsibility attribution; intergroup conflict; Hawk--Dove game

\section{Research Background and Problem}

\subsection{Practical background}

The role of AI in high-stakes decisions is expanding from information analysis and advice to decision agency. In this proposal, AI decision agency has a precise meaning: a human voluntarily transfers final authority to AI, and the AI action directly produces consequences. The human makes the authorization decision, whereas AI makes the final action choice. Because authorization and final action belong to different actors, AI delegation may affect both conflict behavior and subsequent judgments of responsibility.

Intergroup conflict provides a consequential setting for studying this problem. A decision-maker's choice affects not only the decision-maker but also other members of the same group. It also interacts with the opposing group's choice to determine both groups' payoffs. After the outcome, evaluators are not detached observers. They are group members who actually benefited or were harmed. The group design therefore introduces consequences for other real participants, strategic interdependence, and payoff-based evaluation positions.

\subsection{Prior research}

Delegation research shows that transferring authority to an agent may change both behavior and subsequent responsibility attribution. \citet{BartlingFischbacher2012} allowed receivers affected by an unfair allocation to evaluate decision-makers through costly punishment. Delegation shifted much of the punishment, and the responsibility attribution inferred from it, from the delegator to the agent. \citet{KobisEtAl2025} further showed that some AI-delegation interfaces change dishonest behavior with real payoff consequences. This line of research primarily asks whether people delegate and what action follows.

AI-responsibility research mainly asks how people assign responsibility after an outcome has already occurred. \citet{ShankEtAl2019} found that different human--AI decision structures changed the moral fault assigned to humans, whereas AI fault did not necessarily change in parallel. In a shared-control setting, \citet{AwadEtAl2020} rated human drivers and automated systems separately and found that their blame ratings could change asymmetrically rather than being mechanically forced to move in opposite directions. \citet{ArnestadEtAl2024} showed that the mere existence of a manual-control option can increase human blame for AI mistakes. \citet{DongBocian2024} found that self-interest and group relationships affect how people attribute moral responsibility to AI.

Research on AI and conflict indicates that AI involvement may change crisis judgments and escalation choices. \citet{HorowitzLinGreenberg2022} compared judgments about the use of force when intelligence was analyzed by AI or by humans, and compared retaliation judgments after accidents involving an AI-enabled weapon or only human operators. \citet{RiveraEtAl2024} examined escalation behavior by language models in simulated military and diplomatic crises. These studies establish that AI can enter conflict processes, but they do not jointly observe voluntary human delegation, real group payoffs, and subsequent responsibility judgments concerning both the human and AI.

\subsection{Research gap}

The relevant literatures study delegation, conflict behavior, and responsibility attribution separately. Delegation studies usually treat behavior as the main outcome. AI-responsibility studies commonly begin with a completed scenario. AI-conflict studies mainly measure support for force, retaliation judgments, or model escalation choices. Real decision processes, however, unfold in sequence: a human first decides whether to transfer final authority to AI; the human or AI then makes a conflict choice; both sides' actions jointly produce payoff outcomes; and group members subsequently judge responsibility.

Among the representative studies on delegation, AI responsibility, and AI conflict reviewed here, none uses one framework to identify both of the following questions: whether providing an AI decision agent changes final conflict escalation, and whether assigning final authority to the human or AI changes responsibility judgments when the action and outcome are held constant. These studies also do not ask actual benefiting and harmed group members to evaluate their own human and AI. It therefore remains unclear whether people in different payoff positions understand responsibility after AI delegation in the same way.

\section{Research Questions}

The main research question is:

\begin{quote}
In intergroup conflict, does the introduction of an AI decision agent change conflict escalation and group members' responsibility judgments of the human and AI?
\end{quote}

The main question is divided into three specific research questions:

\textbf{RQ1: Delegation.} When an AI decision agent is available, are humans willing to transfer final authority to AI?

\textbf{RQ2: Escalation.} Does making AI delegation available change final conflict escalation?

\textbf{RQ3: Responsibility attribution.} After the conflict outcome, how do benefiting and harmed groups evaluate the responsibility of their own human decision-maker and AI decision agent?

RQ1 records actual delegation within the AI Delegation condition. RQ2 uses randomized availability of AI delegation to identify an overall change in escalation. RQ3 compares independent responsibility ratings for the human and AI and uses a matched-record study to identify the effect of final authority. The decision-maker's ratings of self and AI are reported as a supplementary comparison and are not pooled with the main benefiting-versus-harmed comparison. The three questions belong to the same real-world process, but they do not form a fully identified causal mediation chain.

\section{Theoretical Foundations}

\subsection{Delegation and responsibility shifting}

\term{Delegation} means that a decision-maker voluntarily transfers authority over a choice to an agent. It excludes advice-taking, joint decision-making, and a final human confirmation. In this proposal, delegation occurs at the individual decision-maker level and is an observed behavioral variable. \citet{BartlingFischbacher2012} show that delegation can change subsequent responsibility attribution, while \citet{KobisEtAl2025} show that AI delegation can change behavior with material consequences.

These studies support two competing explanations. Under the first, responsibility follows final authority and shifts toward the agent. Under the second, the human remains responsible for authorizing the transfer even after giving up the final action. The present study evaluates the human and AI separately to distinguish these explanations.

\subsection{Responsibility judgments with multiple actors}

The Triangle Model of Responsibility proposed by \citet{SchlenkerEtAl1994} relates responsibility judgments to connections among prescriptions, identity, and events. The Path Model of Blame developed by \citet{MalleEtAl2014} further distinguishes the roles of causality, intentionality, obligation, and capacity to prevent an outcome. Together, these theories imply that final control is not the sole basis of responsibility judgment.

In AI delegation, the human authorizes the transfer and AI makes the final action choice. Both may therefore become targets of responsibility judgment. This proposal does not combine role obligation, decision control, and causal contribution into a new responsibility construct, nor does it estimate their separate causal effects. Existing responsibility theories instead organize the competing outcomes: responsibility may follow final authority to AI, or it may remain attached to the human because the human authorized the transfer.

\subsection{Evaluator payoff position}

The culpable-control model of \citet{Alicke2000} proposes that an evaluator's general assessment of an actor and outcome can influence the interpretation of control and causal evidence. \citet{DongBocian2024} extend this issue to AI and show that people attribute moral responsibility to AI differently when evaluating themselves or ingroup members. Their studies manipulate self versus other and ingroup versus outgroup relations; they do not directly test benefiting and harmed positions created by real monetary outcomes.

This proposal defines \term{evaluator payoff position} as the position of benefit or harm created by realized experimental payoffs at the time of evaluation. It is an evaluator-level outcome position, not a personality trait and not a synonym for general group identification. In Study 1, benefiting and harmed positions are jointly produced by both groups' choices. In Study 2, the final decision-maker is randomized within each payoff position. The design asks whether the same change in final authority produces different responsibility judgments across payoff positions.

\subsection{Social preferences and the group design}

The inequity-aversion model of \citet{FehrSchmidt1999} allows individual utility to depend on payoff differences between self and others. The intention-based reciprocity model of \citet{Rabin1993} allows behavior to depend on beliefs about another actor's kindness, hostility, and available actions. The present study does not estimate the structural parameters of either model. They instead justify using a group-based conflict task: when a choice affects real participants and produces unequal payoffs, behavior may reflect concern for others' outcomes and expectations about an opponent's likely behavior and intentions, rather than risk preference alone. Because the first round does not reveal the opponent's action before choice, the design does not identify reciprocity in response to an observed action.

\section{Conceptual Model and Hypotheses}

\subsection{Conceptual model}

The research process is summarized as follows:

\begin{equation*}
\underbrace{T}_{\substack{\text{random availability of}\\\text{AI delegation}}}
\longrightarrow
\underbrace{D}_{\substack{\text{actual}\\\text{delegation}}}
\longrightarrow
\underbrace{E}_{\substack{\text{final}\\\text{escalation}}}
\longrightarrow
\underbrace{P}_{\substack{\text{payoff}\\\text{position}}}
\longrightarrow
\underbrace{R}_{\substack{\text{responsibility}\\\text{rating}}}.
\end{equation*}

$T$ is the randomized treatment in Study 1, $D$ is the participant's voluntary delegation choice, $E$ is the final escalation outcome, $P$ is the benefiting or harmed payoff position, and $R$ is the responsibility rating of the human or AI. The sequence describes the order of the study; it does not imply that every arrow is causally identified. Study 1 supports a causal interpretation of $T\rightarrow E$. Study 2 supports a causal interpretation of final decision-maker $F\rightarrow R$. Actual delegation $D$ is not randomized in Study 1.

\subsection{Hypotheses}

\textbf{H1: AI delegation availability and conflict escalation.} Relative to \Baseline{}, \AID{} will change the probability of choosing escalation. The direction depends on actual delegation uptake, the fixed AI decision rule, and possible changes in human choices among those who decline delegation. A directional prediction will be preregistered after the AI rule and pilot results are fixed.

\textbf{H2: Final decision-maker and rated target.} Holding the final action and outcome constant, an AI-made final choice will produce a higher AI responsibility rating than the same choice recorded as having been made by a human.

The joint ratings may display the following patterns. The third would not support H2; these patterns are not exhaustive:

\begin{itemize}
\item \textbf{Responsibility transfer:} AI responsibility increases and human responsibility decreases.
\item \textbf{Shared responsibility:} AI responsibility increases while human responsibility remains high.
\item \textbf{Responsibility not fully taken up:} human responsibility decreases without a clear increase in AI responsibility.
\end{itemize}

These labels describe the joint movement of two responsibility ratings. They are not presented as three established psychological constructs, and they do not assume a fixed total amount of responsibility.

\textbf{H3: Interaction between payoff position and rated target.} Benefiting and harmed groups will form different human--AI responsibility-rating profiles. Existing evidence does not determine which group will necessarily assign more responsibility to AI. This hypothesis concerns $P\times Q$, the difference in rating profiles by payoff position. The separate $F\times P\times Q$ contrast tests whether the effect of the recorded final-decision arrangement on the two ratings differs by position. Their primary or supplementary status will be specified before confirmatory analysis; neither comparison imposes a directional prediction.

\section{Research Design}

\subsection{Study 1: Voluntary AI delegation in group conflict}

Each experimental unit contains two opposing three-person groups. Each group consists of one decision-maker and two members without decision authority. The decision-maker's choice determines the group's action, and all three members share the group's payoff outcome. The two groups choose simultaneously. Neither side can observe the opponent's current action before making its own choice.

Decision-makers are randomly assigned to one of two conditions:

\begin{itemize}
\item \Baseline{}: the human must make the final choice.
\item \AID{}: an AI decision agent is available, and the human may retain final authority or voluntarily implement pure delegation.
\end{itemize}

Pure delegation has four restrictions. The AI and human receive the same task information. Delegation occurs before the final action. Once generated, the AI choice is implemented directly. The human cannot preview, modify, veto, withdraw, or reconfirm the AI choice. At the time of choice, the opponent does not know whether AI delegation is available to the focal group or who makes its final action.

The randomized independent variable in Study 1 is the availability of AI delegation. Actual delegation is chosen by the participant and is treated as the outcome for RQ1 and as a conditional variable in later analysis. Differences between delegators and non-delegators cannot be interpreted as the causal effect of delegation itself.

\subsection{Group-based Hawk--Dove game}

The conflict task is a Hawk--Dove game. Hawk represents escalation and Dove represents restraint. The standard payoff structure is shown in Table~\ref{tab:hawk-dove}, where $V$ is the value of the contested resource and $C$ is the cost of mutual escalation, with $C>V>0$~\citep{MaynardSmithPrice1973}.

\begin{table}[H]
\centering
\caption{Standard Hawk--Dove payoff structure}
\label{tab:hawk-dove}
\begin{tabular}{lcc}
\toprule
\textbf{Focal group choice} & \textbf{Opponent chooses Dove} & \textbf{Opponent chooses Hawk}\\
\midrule
Dove & $(V/2,V/2)$ & $(0,V)$\\
Hawk & $(V,0)$ & $((V-C)/2,(V-C)/2)$\\
\bottomrule
\end{tabular}
\end{table}

The task supplies three features required by the proposal. Unilateral escalation may yield a higher payoff while harming the other group. Mutual escalation imposes a cost on both groups. Finally, the groups' actions jointly determine payoffs, so a decision-maker cannot know the group's realized payoff position before choosing. Exact monetary values and the within-group allocation rule will be fixed after pilot testing and preregistered.

\subsection{Study 1 procedure}

\begin{enumerate}
\item Participants complete informed consent, role instructions, and comprehension checks.
\item Participants are randomly assigned to two three-person groups, and one decision-maker is randomly selected in each group.
\item All participants receive the payoff rules, information boundaries, and decision-authority rules. Group members do not communicate during the decision task.
\item Decision-makers are randomly assigned to C0 or C1.
\item C0 decision-makers choose Hawk or Dove directly. C1 decision-makers either retain control or delegate to AI.
\item Both groups' final actions are implemented simultaneously, and group payoffs are realized.
\item Participants observe their own authorization record, both final actions, and the payoff outcome.
\item In C0, members rate only their human decision-maker. In C1, they rate the human and AI separately; when delegation is declined, the record states that AI did not make the final choice. Decision-makers provide corresponding self-ratings separately. Missing AI ratings in C0 are not coded as zero.
\item Participants complete final comprehension checks and supplementary measures.
\end{enumerate}

The first-round between-subject comparison provides the primary test. If a second round is added, it will be used only to explore whether the previous payoff and AI experience change later delegation intentions or choices. It will not be pooled with the first-round primary result.

\subsection{Study 2: Responsibility ratings for matched records}

Study 2 identifies the effect of the recorded delegation and final-decision arrangement on responsibility ratings. Participants view a completed conflict record presented as concerning their side. The source of these records and whether and how their payoffs affect evaluators' actual bonuses remain to be specified. A described benefiting or harmed role is not, by itself, a real financial position. Until the payment link is implemented, Study 2 claims are restricted to evaluations of assigned records. If participants' own Study 1 records are used, the protocol must also resolve how randomized record information relates to events they already observed. The payoff rule, focal action, opponent action, and focal outcome are held constant. Only the recorded final decision-maker is randomized:

\begin{itemize}
\item the human retained final authority and personally made the choice; or
\item the human voluntarily delegated and AI made the same final choice.
\end{itemize}

Both versions state that AI delegation had been available, so participants rate both the human decision-maker and the AI decision agent in every condition. Matched sets are constructed separately for benefiting and harmed records. The final decision-maker is randomized within each payoff position. Study 2 can therefore test whether the same action and outcome receive different human and AI responsibility ratings when AI rather than a human makes the final choice, and whether that effect differs by payoff position.

\subsection{Measurement of responsibility}

Responsibility refers to the evaluator's judgment of how much each actor should be held responsible for the group's final conflict choice. It is a judgment rating, not a measure of legal liability, compensation, organizational sanction, or experimental punishment.

In Study 1, C0 participants receive only the human item; C1 participants receive both items. Both Study 2 conditions describe AI delegation as available and include both items. The questions are:

\begin{quote}
For the final conflict choice made by your group, how much responsibility should the group's human decision-maker bear?
\end{quote}

\begin{quote}
For the final conflict choice made by your group, how much responsibility should the group's AI decision agent bear?
\end{quote}

Both ratings range from 0 (no responsibility at all) to 100 (full responsibility), and question order is randomized. The ratings are independent and are not constrained to sum to 100. \citet{AwadEtAl2020} separately assessed blame assigned to humans and automated systems in shared control. \citet{ShankEtAl2019} separately assessed the moral fault of humans and AI, while \citet{DongBocian2024} measured moral responsibility attributed to AI. These studies support treating the human and AI as separate targets of evaluation. The present responsibility items are not claimed to replicate any one of these existing measures exactly.

Payoff position is produced by the group's shared outcome and recorded for each evaluator; members of one group are not independent assignments to a payoff condition. The reference payoff, records included in the primary benefit--harm comparison, and treatment of mutual escalation or mutual restraint will be fixed before confirmatory analysis. These symmetric outcomes are reported separately rather than pooled into an undefined ``joint outcome.''

\subsection{Core variables}

Table~\ref{tab:core-variables} summarizes the variables and their interpretation.

\begin{table}[H]
\centering
\small
\caption{Core variables and their roles}
\label{tab:core-variables}
\begin{tabularx}{\textwidth}{@{}>{\RaggedRight\arraybackslash}p{0.18\textwidth}>{\RaggedRight\arraybackslash}p{0.22\textwidth}YY@{}}
\toprule
\textbf{Variable role} & \textbf{Variable} & \textbf{Operationalization or measure} & \textbf{Interpretive scope}\\
\midrule
Randomized independent variable & AI delegation availability, $T$ & C0 versus C1 & Study 1 treatment\\
Behavioral outcome & Actual delegation, $D$ & Human control versus AI delegation within C1 & Voluntary choice; no independent causal interpretation\\
Primary behavioral dependent variable & Final escalation, $E$ & Dove versus Hawk & RQ2\\
Outcome position & Payoff position, $P$ & Benefit/harm records; symmetric outcomes reported separately & Jointly produced by both groups' actions\\
Randomized independent variable & Final decision-maker, $F$ & Human versus AI & Study 2 treatment\\
Repeated-measures factor & Rated target, $Q$ & Human versus AI & Distinguishes the two responsibility targets\\
Primary attribution dependent variable & Responsibility rating, $R$ & Independent 0--100 rating & RQ3\\
Supplementary comparison & Decision-maker self-rating & Decision-maker rates self and AI & Analyzed separately from members' external ratings\\
\bottomrule
\end{tabularx}
\end{table}

\section{Analysis and Causal Identification}

\subsection{Study 1: Delegation and escalation}

Actual delegation is reported only within C1, including its overall rate, confidence interval, and associations with preregistered measures. Because delegation is a voluntary choice, comparisons between delegators and non-delegators receive descriptive rather than causal interpretation.

The primary Study 1 model is given in Equation~\ref{eq:study_1}:

\begin{equation}
E_{gs}=\alpha_1+\beta_TT_{gs}+u_s+\varepsilon_{gs}.
\label{eq:study_1}
\end{equation}

$E_{gs}$ is the final escalation outcome of group $g$ in experimental session $s$, $T_{gs}$ indicates whether AI delegation is available, and $u_s$ captures session-level variation. The main parameter, $\beta_T$, is the intention-to-treat effect of making AI delegation available on final escalation. It is not the causal effect of actual delegation.

Because both groups' actions jointly determine the conflict outcome, standard errors will be clustered at least at the conflict-pair or experimental-session level. Mutual escalation, mutual restraint, and asymmetric action profiles are reported as pair-level supplementary outcomes.

Study 1 also compares human responsibility ratings between C0 and C1. Within C1, joint human and AI ratings are described in relation to actual delegation, group payoff position, and rated target. These conditional associations do not identify the independent causal effects of delegation or payoff position. Evaluator, group, and session dependence will be represented in the analysis; absent C0 AI ratings are never replaced by zeros.

\subsection{Study 2: Responsibility ratings}

The primary Study 2 model is given in Equation~\ref{eq:study_2}:

\begin{equation}
R_{igt}=\alpha_2+\boldsymbol{\beta}^{\mathsf{T}}\mathbf{X}_{igt}+u_i+\varepsilon_{igt}.
\label{eq:study_2}
\end{equation}

$R_{igt}$ is evaluator $i$'s responsibility rating of target $t$ in record $g$. The vector $\mathbf{X}_{igt}$ contains final decision-maker $F$, payoff position $P$, rated target $Q$, all two-way interactions, and the three-way interaction. The $F\times Q$ interaction tests whether the effect of the recorded arrangement differs between human and AI ratings. Its significance alone does not establish that one rating rises while the other falls; joint marginal means determine that pattern. The $F\times P\times Q$ interaction tests whether this shift differs between benefiting and harmed positions.

Study 2 randomizes the final decision-maker within each payoff position. The effect of the randomized record arrangement can therefore receive a causal interpretation. This arrangement jointly changes whether the human delegated and who made the final choice; it does not isolate decision control as a separate psychological mechanism. The contrast between benefiting and harmed records remains a position comparison because the two positions correspond to different payoff outcomes.

\subsection{Joint interpretation of responsibility ratings}

Responsibility transfer, shared responsibility, and responsibility not fully taken up are classified from the joint movement of the human and AI ratings. The analysis will use preregistered marginal means, confidence intervals, and a smallest effect size of interest. Claims of no meaningful change will rely on equivalence tests or a preregistered smallest effect size rather than a nonsignificant $p$-value. The analysis compares the two ratings jointly and does not compare their simple sum across conditions. The ratings do not measure an objective total amount of responsibility.

\section{Expected Contributions}

First, the proposal places AI delegation, conflict escalation, and post-outcome responsibility judgments within the same group-conflict process. Study 1 identifies the overall effect of making AI delegation available on escalation, while Study 2 identifies the effect of the recorded delegation and final-decision arrangement on responsibility ratings. The two studies provide separate evidence about behavioral and attribution outcomes.

Second, the proposal measures human and AI responsibility independently. The design can distinguish responsibility transfer, shared responsibility, and responsibility not fully taken up without mechanically forcing one rating to rise when the other falls. The joint ratings test whether transferring final authority is accompanied by lower human responsibility or continued responsibility for both actors. Continued human responsibility is consistent with responsibility for authorization, but does not by itself identify that mechanism.

Third, Study 1 collects responsibility judgments from members whose groups receive real payoff consequences. Study 2 tests whether the effect of the recorded decision arrangement differs between benefiting and harmed records. Extending the latter result to actual financial positions requires a specified and implemented link between the record and the evaluator's bonus.

These contributions have clear limits. The absence of a human-agent condition means that the study cannot establish an effect unique to AI rather than delegation more generally. Actual delegation in Study 1 is self-selected, so differences between delegators and non-delegators cannot be interpreted as the independent causal effect of delegation. The two studies also do not establish that post-outcome responsibility judgments caused earlier delegation or escalation.

\section{Ethics, Pilot Testing, and Preregistration}

Ethics approval or an exemption determination will be obtained before formal recruitment. Base compensation and performance bonuses will be separated, and the lowest experimental outcome will not produce a negative payment. Any use of pregenerated AI output, simulated opponents, or stored decision records will be disclosed in the ethics application and debriefing.

Pilot testing will assess comprehension of the payoff matrix, AI-delegation uptake, the escalation tendency of the fixed AI decision rule, the frequency of asymmetric payoff outcomes, and comprehension of the responsibility items. Study 1 and Study 2 will receive separate simulation-based power analyses. Preregistration will fix the AI action rule, primary treatments, primary outcomes, clustering level, responsibility items, criteria for classifying responsibility patterns, and multiple-testing procedures.




\begingroup
\RaggedRight
\begin{thebibliography}{99}
\RaggedRight
\interlinepenalty=10000

\bibitem[Alicke(2000)]{Alicke2000}
Alicke, M. D. (2000). Culpable control and the psychology of blame. \textit{Psychological Bulletin, 126}(4), 556--574. \url{https://doi.org/10.1037/0033-2909.126.4.556}

\bibitem[Arnestad et al.(2024)]{ArnestadEtAl2024}
Arnestad, M. N., Meyers, S., Gray, K., et al. (2024). The existence of manual mode increases human blame for AI mistakes. \textit{Cognition, 252}, 105931. \url{https://doi.org/10.1016/j.cognition.2024.105931}

\bibitem[Awad et al.(2020)]{AwadEtAl2020}
Awad, E., Levine, S., Kleiman-Weiner, M., et al. (2020). Drivers are blamed more than their automated cars when both make mistakes. \textit{Nature Human Behaviour, 4}(2), 134--143. \url{https://doi.org/10.1038/s41562-019-0762-8}

\bibitem[Bartling and Fischbacher(2012)]{BartlingFischbacher2012}
Bartling, B., \& Fischbacher, U. (2012). Shifting the blame: On delegation and responsibility. \textit{The Review of Economic Studies, 79}(1), 67--87. \url{https://doi.org/10.1093/restud/rdr023}

\bibitem[Dong and Bocian(2024)]{DongBocian2024}
Dong, M., \& Bocian, K. (2024). Responsibility gaps and self-interest bias: People attribute moral responsibility to AI for their own but not others' transgressions. \textit{Journal of Experimental Social Psychology, 111}, 104584. \url{https://doi.org/10.1016/j.jesp.2023.104584}

\bibitem[Fehr and Schmidt(1999)]{FehrSchmidt1999}
Fehr, E., \& Schmidt, K. M. (1999). A theory of fairness, competition, and cooperation. \textit{The Quarterly Journal of Economics, 114}(3), 817--868. \url{https://doi.org/10.1162/003355399556151}

\bibitem[Horowitz and Lin-Greenberg(2022)]{HorowitzLinGreenberg2022}
Horowitz, M. C., \& Lin-Greenberg, E. (2022). Algorithms and influence: Artificial intelligence and crisis decision-making. \textit{International Studies Quarterly, 66}(4), sqac069. \url{https://doi.org/10.1093/isq/sqac069}

\bibitem[K\"{o}bis et al.(2025)]{KobisEtAl2025}
K\"{o}bis, N., Rahwan, Z., Rilla, R., et al. (2025). Delegation to artificial intelligence can increase dishonest behaviour. \textit{Nature, 646}(8083), 126--134. \url{https://doi.org/10.1038/s41586-025-09505-x}

\bibitem[Malle et al.(2014)]{MalleEtAl2014}
Malle, B. F., Guglielmo, S., \& Monroe, A. E. (2014). A theory of blame. \textit{Psychological Inquiry, 25}(2), 147--186. \url{https://doi.org/10.1080/1047840X.2014.877340}

\bibitem[Maynard Smith and Price(1973)]{MaynardSmithPrice1973}
Maynard Smith, J., \& Price, G. R. (1973). The logic of animal conflict. \textit{Nature, 246}, 15--18. \url{https://doi.org/10.1038/246015a0}

\bibitem[Rabin(1993)]{Rabin1993}
Rabin, M. (1993). Incorporating fairness into game theory and economics. \textit{American Economic Review, 83}(5), 1281--1302.

\bibitem[Rivera et al.(2024)]{RiveraEtAl2024}
Rivera, J.-P., Mukobi, G., Reuel, A., et al. (2024). Escalation risks from language models in military and diplomatic decision-making. In \textit{Proceedings of the 2024 ACM Conference on Fairness, Accountability, and Transparency} (pp. 836--898). \url{https://doi.org/10.1145/3630106.3658942}

\bibitem[Schlenker et al.(1994)]{SchlenkerEtAl1994}
Schlenker, B. R., Britt, T. W., Pennington, J., et al. (1994). The triangle model of responsibility. \textit{Psychological Review, 101}(4), 632--652. \url{https://doi.org/10.1037/0033-295X.101.4.632}

\bibitem[Shank et al.(2019)]{ShankEtAl2019}
Shank, D. B., DeSanti, A., \& Maninger, T. (2019). When are artificial intelligence versus human agents faulted for wrongdoing? Moral attributions after individual and joint decisions. \textit{Information, Communication \& Society, 22}(5), 648--663. \url{https://doi.org/10.1080/1369118X.2019.1568515}

\end{thebibliography}
\endgroup

\end{document}
